
\subsection{Ergebnisse}
Diese Arbeit betrachtet \acp{CbCC} in dem Kontext von
\acp{TEA} losgelöst von den tatsächlichen Angriffen.
Dadurch wird allgemeiner das Problem der Laufzeit der
Übertragung von Daten während der Transient Execution
analysiert.

Es wird das erste Mal die praktische Umsetzung
der bitweisen Encodierung, im Kontext von \acp{TEA},
demonstriert und mit der Array-Index-Encodierung verglichen.
Dabei zeigt sich, dass die bitweise Encodierung, bezüglich
der Anzahl der encodierten Bits, begrenzt ist.
Die Array-Index-Encodierung weist dagegen ein exponentiell
steigendes Laufzeitverhalten bei steigender Bitanzahl auf.
Aus diesem Grund wird der Durchsatz abhängig von der Bitanzahl
und verschiedenen messbaren Parametern modelliert und eine
Funktion für die Berechnung der optimalen Encodierung
hergeleitet.

Des Weiteren wird die obere Grenze
$\frac{1}{T_P+T_R}$ für den Durchsatz von \acp{CbCC} im Kontext
von Transient Execution ermittelt. Die ermittelte
Durchsatzfunktion $D(n_a)$ kann zudem auch als obere Grenze für
den Durchsatz eines \ac{CbCC} interpretiert werden, insofern für
$T_R$ die geringere Laufzeit zwischen einer gecachten und einer
nicht gecachten Adresse verwendet wird.
Zuletzt wird an einem bereits bekannten \ac{POC} verdeutlicht, wie
stark ein praktischer \ac{TEA} durch die Optimierung des
verwendeten \ac{CbCC} und damit das ausgehende Risiko verändert
werden kann.

Auch wenn Gadgets behoben und Mitigationen gegen
\acp{TEA} in Mikrocode entwickelt werden, bleibt bei
bestehender Hardware das zugrunde liegende Problem der
Schwachstellen bestehen (\cite[3,10]{hertogh_rain_nodate}).
Neue \acp{CPU} verwenden außerdem weiter spekulative
Ausführung, um die Performance weiter zu steigern. Dabei
handelt es sich nicht nur um Branch-Prediction, sondern auch
um neue Predictoren (\cite[5]{kim_slap_2025}). Diese neuen
Performanceverbesserungen bieten auch neue Angriffsflächen
für \acp{TEA}. Diese Art von Angriffen wird also in Zukunft
bestehen bleiben und kann präventiv nur begrenzt werden.
Aus diesem Grund ist besonders für Endgeräte die Laufzeit
eines möglichen Angriffs relevant. Die hergeleitete Funktion
für den Durchsatz $D(n_a)$ ermöglicht die Einschätzung
der Laufzeit dieser Angriffe unabhängig von der ausgenutzten
Schwachstelle.

\subsection{Ausblick}
Die theoretische Betrachtung des Durchsatzes beachtet nicht
die Fehlerrate des \ac{CbCC}. Besonders in Bezug auf die
beiden Encodierungen müsste für einen realistischen Angriff
und dessen Limits auch die Korrektheit der übertragenen Daten
und damit auch Fehlerkorrekturverfahren in Betracht gezogen
werden. Wie die Sektion \ref{paragraph:error-probability} zeigt,
verhält sich die Fehlerwahrscheinlichkeit der beiden Encodierungen
aufgrund der gegebenen Eigenschaften sehr unterschiedlich.
Dies erschwert die Betrachtung der Fehlerwahrscheinlichkeiten
für die verallgemeinerte Encodierung deutlich.

Eine weitere offene Frage ist, wie Cache-Zustände optimal während
der Transient Execution geändert werden können. Bekannt ist
beispielsweise nicht, ob andere Instruktionen, wie
\verb|clflush|, in demselben transient Fenster häufiger
ausgeführt werden könnten als die \verb|prefetch0|-Instruktion.
Des Weiteren könnte die Anzahl der Instruktionen für das
Encodieren optimiert werden, um möglichst viele Zustände zu
ändern.

Die aktuelle Berechnung für den Durchsatz $D(n_a)$ kann
vermutlich durch das Hinzufügen weiterer Parameter
präzisiert werden. Dafür könnte beispielsweise ein Parameter
für die Zeit zwischen zwei Übertragungen hinzugefügt oder
die unterschiedliche Laufzeit zwischen einer gecachten
und einer nicht gecachten Adresse beachtet werden. Dadurch
könnte der maximale Durchsatz weiter eingegrenzt werden.
Letztendlich sollten mathematische Modelle empirisch
verifiziert werden, wofür die Bearbeitungszeit für diese
Arbeit nicht ausgereicht hat. An einem entsprechenden
Programm zur Ermittlung der nötigen Parameter und zum Messen
des Durchsatzes für gegebene $n_b$ und $n_a$ habe ich bereits
zum Zeitpunkt der Abgabe dieser Arbeit gearbeitet. Die
Entwicklung des Programms setze ich fort, mit dem Ziel,
empirische Werte des Durchsatzes mit der theoretischen Grenze
sowie der Durchsatzfunktion zu vergleichen.
